 \documentclass{article}

  \usepackage{fancyhdr}
  \usepackage{amsmath}
  \usepackage{amsthm}
  \usepackage{amsfonts}
  \usepackage{tikz}
  \usepackage[plain]{algorithm}
  \usepackage{algpseudocode}

  \usetikzlibrary{automata,positioning}

  %
  % Basic Document Settings
  %

  \topmargin=-0.45in
  \evensidemargin=0in
  \oddsidemargin=0in
  \textwidth=6.5in
  \textheight=9.0in
  \headsep=0.25in

  \linespread{1.1}

  \pagestyle{fancy}
  \lhead{\hmwkAuthorName}
  \chead{} % empty
  \rhead{\hmwkHeaderTitle}
  \cfoot{\thepage}

  \renewcommand\headrulewidth{0.4pt}
  \renewcommand\footrulewidth{0pt}

  \setlength\parindent{0pt}

  %
  % Homework Details
  %

  \newcommand{\hmwkTitle}{Homework\ 1}
  \newcommand{\hmwkDueDate}{Tuesday, January 13, 2026,\ 11{:}59 pm (Thai time)}
  \newcommand{\hmwkClass}{Data Structure\ and Abstractions}
  \newcommand{\hmwkClassTime}{Section 2}
  \newcommand{\hmwkClassInstructor}{Kanat Tangwongsan}
  \newcommand{\hmwkAuthorName}{\textbf{Jiraroj Wiruchpongsanon (6781617)}}

  % What appears in the top-right header
  \newcommand{\hmwkHeaderTitle}{Data Structure and Abstractions: Homework 1}

  %
  % Title Page
  %

  \title{
      \vspace{2in}
      \textmd{\textbf{\hmwkClass:\ \hmwkTitle}}\\
      \normalsize\vspace{0.1in}\small{Due\ on\ \hmwkDueDate}\\
      \vspace{0.1in}\large{\textit{\hmwkClassInstructor\ \hmwkClassTime}}
      \vspace{3in}
  }

  \author{\hmwkAuthorName}
  \date{}

  \renewcommand{\part}[1]{\textbf{\large Part \Alph{partCounter}}\stepcounter{partCounter}\\}

  %
  % Helper Commands
  %

  \newcounter{partCounter}
  \newcommand{\solution}{\textbf{\large Solution}}

  \newtheorem*{theorem}{Theorem}

  \theoremstyle{remark}
  \newtheorem*{remark}{Remark}

  \begin{document}

  \maketitle
  \newpage

  \section*{Important Notes from the Assignment}

  \textit{Typesetting Requirement.} Written answers must be typed and compiled into
  \texttt{hw1.pdf}; scans of handwriting are not accepted.

  \textit{AI Usage.} You must write every element of the code yourself. Generative AI may assist
  with syntax questions and pointers only.

  \textit{Hand-In.} Package \texttt{hw1.pdf} together with the required Java source files in the
  top-level directory of \texttt{a1.zip} before the deadline. Filenames are case sensitive.

  \textit{Collaboration.} You may collaborate conceptually but must produce your own write-up and
  code; always list collaborators and sources inside the submission.

  \newpage

  \section*{Problem 1 (4 points)}

  \textbf{Task 2 --- Oh My Magical Math (OHMM\ldots{})}

  The first $n$ even numbers are $0, 2, 4, \ldots, 2(n-1)$. This task involves computing the sum
  of the first
  $n$ even numbers and understanding their performance implications. In addition to the Java
  code in
  \texttt{SumEven.java}, answer the following in \texttt{hw1.pdf}. Show all work.

  \begin{enumerate}
      \item[(a)] Using the method discussed in Lecture 1 (Gauss's pairing trick), derive a simple
  formula for the sum of the first $n$ even numbers. Derive it directly without referencing other
  formulas.
      \item[(b)] If we replace the loop in \texttt{sumEven} with this closed-form formula,
  approximately how many steps does the implementation need as a function of $n$?
  \end{enumerate}

  \medskip
  \solution

  \begin{enumerate}
      \item[(a)] let $L$ be the sum of the first n even numbers, $L$ can be written as
        \begin{equation}\tag{a}\label{eq:L-def}
          L = 0 + 2 + 4 + 6 + ... + 2(n-2) + 2(n-1)
        \end{equation}

          or, with the commutative property of real numbers,
        \begin{equation}\tag{a'}
            L = 2(n-1) + 2(n-2) + 2(n-3) + 2(n-4) + ... + 2 + 0
        \end{equation}

          pair up the $k^{th}$ term of $a$ with the $k^{th}$ term of $a'$, when $k = 1, 2, 3, ..., n$, we'll get
        \[
          \begin{aligned}
          2L &= 2(n-1) + (2 + 2(n-2)) + (4 + 2(n-3)) + (6 + 2(n-4)) + ... + (2(n-2) + 2) + 2(n-1) \\
             &= \underbrace{2(n-1) + 2(n-1) + ... + 2(n-1)}_{n}
          \end{aligned}
        \]

        we'll get 
          \[2L = n \cdot 2(n-1)\]

          which simplifies to $\boxed{L = n(n-1)}$ --- and that's the formular for the sum of the first n even numbers with Gauss's pairing trick. 

      \item[(b)] This 'formula' would have a time complexity of O(1), so roughly 1 step. Though it'd depends on the machine's archetechture, whether it could do multiplication in a single clock cycle, or more --- that'd be 'more precisely' at least 2 steps (for subtraction and multiplication).

  \end{enumerate}

  \end{document}
